\documentclass{beamer}

\usetheme{Madrid}
\usecolortheme{seahorse}
\usepackage[utf8]{inputenc}
\usepackage[T1]{fontenc}
\usepackage{listings}
\usepackage{xcolor}
\usepackage{amsmath}
\usepackage{booktabs}

% --- Configuracion del estilo del codigo C# ---
\lstdefinestyle{CSharpStyle}{
  language=[Sharp]C,
  basicstyle=\scriptsize\ttfamily,
  keywordstyle=\color{blue}\bfseries,
  stringstyle=\color{teal},
  commentstyle=\color{gray}\itshape,
  showstringspaces=false,
  breaklines=true,
  frame=single,
  rulecolor=\color{gray!30},
  tabsize=2,
  captionpos=b,
  numbers=left,
  numberstyle=\tiny\color{gray},
  morekeywords={GameObject,Vector3,Quaternion,Instantiate,Sprite,SpriteRenderer,TextAsset,File,Path,Debug,List,StringSplitOptions}
}

% --- Titulo ---
\title[Arreglos, Matrices y Archivos en C\#]{Arreglos, Matrices y Lectura de Archivos de Texto en C\# para Unity}
\subtitle{Estructuras de Datos y Generaci\'on de Mapas Tipo Cuadr\'icula (Grid)}
\author{Luis Antonio Molina Yampa}
\institute{COM200 - Programacion Avanzada}
\date{\today}

\begin{document}

% --- Portada ---
\begin{frame}
  \titlepage
\end{frame}

% --- Introduccion y Objetivos ---
\begin{frame}{Contenido y Objetivos}
  En esta presentaci\'on se exploran las bases fundamentales del manejo de estructuras de datos en C\# aplicadas al desarrollo de videojuegos en Unity:

  \vspace{0.3cm}
  \begin{itemize}
    \item \textbf{Arreglos unidimensionales ($T[]$):} Declaraci\'on, indexaci\'on y recorridos.
    \item \textbf{Matrices en C\#:} Diferencias entre matrices rectangulares (\texttt{int[,]}) y escalonadas o \textit{jagged} (\texttt{int[][]}).
    \item \textbf{Lectura de archivos \texttt{.txt}:} Uso de \texttt{TextAsset} vs. \texttt{System.IO.File}.
    \item \textbf{Pipeline de parseo:} De texto sin formato a matrices bidimensionales.
    \item \textbf{Generaci\'on de cuadr\'iculas (Grid):} Sistema de coordenadas espaciales estilo Battle City e instanciaci\'on de prefabs en escena.
  \end{itemize}
\end{frame}

% --- Arreglos Unidimensionales: Concepto ---
\begin{frame}{Arreglos Unidimensionales ($T[]$) --- Concepto}
  Un \textbf{arreglo (array)} en C\# es una colecci\'on contigua y de tama\~no fijo que almacena elementos del mismo tipo de dato.

  \vspace{0.3cm}
  \textbf{Caracter\'isticas principales:}
  \begin{itemize}
    \item \textbf{Indexaci\'on en base cero:} Los \'indices van desde \texttt{0} hasta $\text{Length} - 1$.
    \item \textbf{Tama\~no est\'atico:} Una vez instanciado, su capacidad no puede alterarse.
    \item \textbf{Acceso r\'apido en $O(1)$:} Acceso inmediato a cualquier posici\'on mediante su \'indice en memoria.
  \end{itemize}

  \vspace{0.3cm}
  \textbf{Aplicaciones t\'ipicas en Unity:}
  \begin{itemize}
    \item Lista fija de prefabs o sprites de bloques (ladrillo, acero, agua).
    \item Puntos de patrullaje (\textit{waypoints}) predeterminados.
    \item Filas individuales le\'idas de un mapa.
  \end{itemize}
\end{frame}

% --- Declaracion e Inicializacion de Arrays ---
\begin{frame}[fragile]{Declaraci\'on e Inicializaci\'on de Arreglos}
\lstset{style=CSharpStyle}
C\# ofrece tres alternativas comunes para declarar e inicializar arreglos:

\begin{lstlisting}
// 1. Con tamano fijo y valores por defecto (0 para int)
int[] numbers = new int[5]; // { 0, 0, 0, 0, 0 }

// 2. Declaracion con valores explicitos
int[] reglon1 = new int[] { 1, 1, 1, 0, 2, 2, 3 };

// 3. Sintaxis simplificada (el compilador infiere tipo y tamano)
int[] reglon2 = { 1, 0, 0, 0, 0, 0, 1 };
\end{lstlisting}

\vspace{0.2cm}
\textbf{Nota:} Si se crea con \texttt{new int[N]}, todos los elementos se inicializan autom\'aticamente en su valor por defecto (\texttt{0} para n\'umeros, \texttt{null} para objetos).
\end{frame}

% --- Acceso y Modificacion ---
\begin{frame}[fragile]{Acceso, Modificaci\'on y Seguridad}
\lstset{style=CSharpStyle}
El acceso y alteraci\'on de valores se realiza mediante corchetes \texttt{[]}:

\begin{lstlisting}
int[] scores = { 10, 20, 30, 40 };

// Lectura de un elemento
int primerValor = scores[0]; // 10

// Modificacion de un elemento
scores[1] = 99; // El arreglo ahora es { 10, 99, 30, 40 }

// Cantidad total de elementos
int total = scores.Length; // 4
\end{lstlisting}

\vspace{0.2cm}
\begin{alertblock}{Precauci\'on: L\'imites del Arreglo}
Intentar acceder a un \'indice $< 0$ o $\ge \text{Length}$ produce la excepci\'on \texttt{IndexOutOfRangeException}. Siempre debe verificarse: $0 \le i < \text{array.Length}$.
\end{alertblock}
\end{frame}

% --- Recorrido de Arreglos ---
\begin{frame}[fragile]{Recorrido de Arreglos: \texttt{for} vs. \texttt{foreach}}
\lstset{style=CSharpStyle}
\textbf{Bucle \texttt{for}:} Permite conocer y manipular el \'indice actual, adem\'as de modificar elementos:
\begin{lstlisting}
for (int i = 0; i < scores.Length; i++)
{
    Debug.Log($"Elemento en indice {i}: {scores[i]}");
}
\end{lstlisting}

\vspace{0.2cm}
\textbf{Bucle \texttt{foreach}:} Sintaxis m\'as limpia y declarativa; \'util cuando s\'olo se requiere lectura:
\begin{lstlisting}
foreach (int score in scores)
{
    Debug.Log($"Puntaje: {score}");
}
\end{lstlisting}
\end{frame}

% --- Matrices en C#: Comparacion ---
\begin{frame}{Matrices en C\#: Rectangulares vs. Escalonadas}
  C\# admite dos formas nativas para trabajar con dos dimensiones:

  \vspace{0.3cm}
  \begin{table}
    \centering
    \footnotesize
    \begin{tabular}{|l|p{4.2cm}|p{4.2cm}|}
      \hline
      \textbf{Caracter\'istica} & \textbf{Rectangular (\texttt{int[,]})} & \textbf{Escalonada (\texttt{int[][]})} \\ \hline
      \textbf{Estructura} & Bloque continuo fijo & Arreglo de arreglos (\textit{Jagged}) \\ \hline
      \textbf{Longitud filas} & Todas tienen la misma longitud & Cada fila puede tener longitud distinta \\ \hline
      \textbf{Sintaxis acceso} & \texttt{matriz[r, c]} & \texttt{matriz[r][c]} \\ \hline
      \textbf{Consulta tama\~no} & \texttt{GetLength(0)}, \texttt{GetLength(1)} & \texttt{Length}, \texttt{matriz[r].Length} \\ \hline
      \textbf{Parseo archivos} & R\'igido; requiere conocer dimensiones previas & Muy natural para construcci\'on l\'inea por l\'inea \\ \hline
    \end{tabular}
  \end{table}
\end{frame}

% --- Matriz Rectangular ---
\begin{frame}[fragile]{Matriz Rectangular (\texttt{int[,]})}
\lstset{style=CSharpStyle}
Se utiliza cuando la cuadr\'icula tiene dimensiones uniformes y perfectamente definidas:

\begin{lstlisting}
// Declaracion e instanciacion fija (3 filas x 5 columnas)
int[,] mapaFijo = new int[,]
{
    { 1, 1, 1, 0, 2 },
    { 1, 0, 0, 0, 1 },
    { 1, 1, 1, 0, 2 }
};

// Acceso y asignacion directa con coma
int tile = mapaFijo[0, 2]; // fila 0, columna 2 -> 1
mapaFijo[1, 1] = 3;        // asigna 3 en (1, 1)

// Obtencion de dimensiones
int filas = mapaFijo.GetLength(0);    // 3
int columnas = mapaFijo.GetLength(1); // 5
\end{lstlisting}
\end{frame}

% --- Recorrido Matriz Rectangular ---
\begin{frame}[fragile]{Recorrido de Matriz Rectangular}
\lstset{style=CSharpStyle}
Para iterar sobre una matriz rectangular se utiliza un doble bucle anidado utilizando \texttt{GetLength()}:

\begin{lstlisting}
int rows = mapaFijo.GetLength(0);
int cols = mapaFijo.GetLength(1);

for (int r = 0; r < rows; r++)
{
    for (int c = 0; c < cols; c++)
    {
        int cellValue = mapaFijo[r, c];
        // Procesamiento de cada celda (r, c)
    }
}
\end{lstlisting}

\vspace{0.2cm}
\textbf{Observaci\'on:} Todas las filas tienen obligatoriamente la misma cantidad de columnas (\texttt{cols}).
\end{frame}

% --- Matriz Escalonada (Jagged) ---
\begin{frame}[fragile]{Matriz Escalonada (\texttt{int[][]}) --- Arreglo de Arreglos}
\lstset{style=CSharpStyle}
Un arreglo escalonado (\textit{jagged array}) almacena arreglos internos independientes:

\begin{lstlisting}
// Declaracion de un arreglo que contendra 4 filas
int[][] matriz = new int[4][];

// Inicializacion independiente por fila
matriz[0] = new int[] { 1, 1, 1, 0, 2 };
matriz[1] = new int[] { 1, 0, 0, 0, 1 };
matriz[2] = new int[] { 1, 0, 0, 0, 1 };
matriz[3] = new int[] { 1, 1, 1, 0, 2 };

// Acceso mediante corchetes dobles
int valor = matriz[0][2]; // fila 0, columna 2
\end{lstlisting}

\vspace{0.2cm}
\textbf{Ventaja clave:} Permite que cada fila provenga de una cadena o l\'inea le\'ida directamente de un archivo de texto.
\end{frame}

% --- Recorrido Matriz Escalonada ---
\begin{frame}[fragile]{Recorrido de Matriz Escalonada}
\lstset{style=CSharpStyle}
Cada fila consulta su propia longitud con \texttt{matriz[r].Length}:

\begin{lstlisting}
for (int r = 0; r < matriz.Length; r++)
{
    for (int c = 0; c < matriz[r].Length; c++)
    {
        int cell = matriz[r][c];
        // Procesamiento del elemento
    }
}
\end{lstlisting}

\vspace{0.3cm}
\textbf{Composici\'on din\'amica desde arreglos 1D:}
\begin{lstlisting}
int[] filaA = { 1, 1, 1, 0 };
int[] filaB = { 1, 0, 0, 1 };
int[][] mapaV2 = { filaA, filaB };
\end{lstlisting}
\end{frame}

% --- Lectura de Archivos de Texto en Unity ---
\begin{frame}{Lectura de Archivos \texttt{.txt} en Unity}
  Existen dos m\'etodos principales para cargar mapas de texto en Unity:

  \vspace{0.3cm}
  \begin{block}{1. \texttt{UnityEngine.TextAsset} (Recomendado para videojuegos)}
    \begin{itemize}
      \item Unity reconoce autom\'aticamente archivos \texttt{.txt}, \texttt{.json}, \texttt{.csv}.
      \item Se empaqueta en la \textit{build} final del juego.
      \item Multiplataforma nativo (Windows, Mac, Android, iOS, WebGL).
      \item Asignaci\'on visual v\'ia Inspector por \textit{drag and drop}.
    \end{itemize}
  \end{block}

  \begin{block}{2. \texttt{System.IO.File} (Archivos externos / Modding)}
    \begin{itemize}
      \item Lee directamente del disco (\texttt{Application.streamingAssetsPath}).
      \item Permite a los jugadores crear o modificar niveles sin recompilar.
      \item En Android/WebGL requiere consideraciones adicionales de acceso.
    \end{itemize}
  \end{block}
\end{frame}

% --- Carga con TextAsset ---
\begin{frame}[fragile]{M\'etodo 1: \texttt{UnityEngine.TextAsset}}
\lstset{style=CSharpStyle}
\begin{lstlisting}
using UnityEngine;

public class MapLoaderTextAsset : MonoBehaviour
{
    [SerializeField] private TextAsset mapFile;
    [SerializeField] private TextAsset[] allMapFiles; // Multiples niveles

    void Start()
    {
        if (mapFile != null)
        {
            // .text entrega el contenido completo del archivo como string
            string rawContent = mapFile.text;
            Debug.Log($"Contenido cargado:\n{rawContent}");
        }
    }
}
\end{lstlisting}

\vspace{0.2cm}
\textbf{Beneficio:} Cero rutas r\'igidas (\textit{hardcoded paths}) y compatibilidad total con cualquier dispositivo.
\end{frame}

% --- Carga con System.IO.File ---
\begin{frame}[fragile]{M\'etodo 2: \texttt{System.IO.File} (StreamingAssets)}
\lstset{style=CSharpStyle}
\begin{lstlisting}
using System.IO;
using UnityEngine;

public class MapLoaderSystemIO : MonoBehaviour
{
    void LoadExternalMap(string fileName)
    {
        string fullPath = Path.Combine(
            Application.streamingAssetsPath, fileName);

        if (File.Exists(fullPath))
        {
            // Lee todas las lineas a la vez en un arreglo de strings
            string[] lines = File.ReadAllLines(fullPath);
        }
        else
        {
            Debug.LogError($"Archivo no encontrado en: {fullPath}");
        }
    }
}
\end{lstlisting}
\end{frame}

% --- Pipeline de Parseo ---
\begin{frame}{Pipeline de Transformaci\'on: De Texto a Matriz 2D}
  \textbf{Etapas secuenciales del procesamiento de mapas:}

  \vspace{0.3cm}
  \begin{enumerate}
    \item \textbf{Texto sin procesar:} Obtener el \texttt{string} completo desde \texttt{TextAsset.text}.
    \item \textbf{Divisi\'on por l\'ineas:} Separar por saltos de l\'inea (\verb|'\n'|, \verb|'\r'|).
    \item \textbf{Limpieza y filtrado:} Aplicar \texttt{Trim()} e ignorar l\'ineas vac\'ias o comentarios (\texttt{\#}, \texttt{//}).
    \item \textbf{Separaci\'on de celdas:} Dividir cada l\'inea por delimitadores (\verb|','|, espacio, \verb|'\t'|).
    \item \textbf{Conversi\'on num\'erica:} Emplear \texttt{int.TryParse()} para evitar ca\'idas por caracteres inv\'alidos.
    \item \textbf{Construcci\'on de matriz:} Almacenar cada fila en \texttt{List<int[]>} y transformar a \texttt{int[][]}.
  \end{enumerate}
\end{frame}

% --- Implementacion del Parser ---
\begin{frame}[fragile]{Implementaci\'on del Parser: \texttt{MapParser}}
\lstset{style=CSharpStyle}
\begin{lstlisting}
using System;
using System.Collections.Generic;

public static class MapParser
{
    public static int[][] ParseTextToMatrix(string textContent)
    {
        if (string.IsNullOrEmpty(textContent)) return new int[0][];

        string[] rawLines = textContent.Split(
            new[] { '\r', '\n' }, StringSplitOptions.RemoveEmptyEntries);
        List<int[]> matrixList = new List<int[]>();

        for (int r = 0; r < rawLines.Length; r++)
        {
            string line = rawLines[r].Trim();
            if (string.IsNullOrEmpty(line) || line.StartsWith("#") || line.StartsWith("//"))
                continue;

            string[] tokens = line.Split(
                new[] { ',', ' ', '\t' }, StringSplitOptions.RemoveEmptyEntries);
            int[] row = new int[tokens.Length];

            for (int c = 0; c < tokens.Length; c++)
                int.TryParse(tokens[c], out row[c]);

            matrixList.Add(row);
        }
        return matrixList.ToArray();
    }
}
\end{lstlisting}
\end{frame}

% --- Coordenadas del Mapa estilo Battle City ---
\begin{frame}{Sistema de Coordenadas de la Cuadr\'icula (Grid)}
  En juegos 2D tipo \textit{Battle City} o laberintos:
  \begin{itemize}
    \item Las \textbf{columnas ($c$)} avanzan de izquierda a derecha en el eje $+X$.
    \item Las \textbf{filas ($r$)} avanzan de arriba hacia abajo en el eje $-Y$.
  \end{itemize}

  \vspace{0.3cm}
  \textbf{F\'ormula de posici\'on en el espacio del mundo:}
  \[
    \vec{P}(r, c) = \begin{pmatrix} X_{\text{origen}} + c \times \text{cellWidth} \\ Y_{\text{origen}} - r \times \text{cellHeight} \\ 0 \end{pmatrix}
  \]

  \vspace{0.2cm}
  \begin{itemize}
    \item $X_{\text{origen}}, Y_{\text{origen}}$: Posici\'on de la esquina superior izquierda.
    \item $\text{cellWidth}, \text{cellHeight}$: Ancho y alto de cada celda en unidades de Unity.
    \item El signo negativo en $Y$ garantiza que la fila 0 est\'e en la parte superior.
  \end{itemize}
\end{frame}

% --- Instanciacion en Unity: BuildGrid ---
\begin{frame}[fragile]{Instanciaci\'on en Escena: \texttt{BuildGrid}}
\lstset{style=CSharpStyle}
\begin{lstlisting}
public void BuildGrid(int[][] matrix, GameObject wallPrefab, 
                      Sprite[] wallSprites, float cellWidth, float cellHeight)
{
    Vector3 origin = transform.position;

    for (int r = 0; r < matrix.Length; r++)
    {
        for (int c = 0; c < matrix[r].Length; c++)
        {
            int tileType = matrix[r][c];
            if (tileType == 0) continue; // 0 = celda vacia

            Vector3 worldPos = new Vector3(
                origin.x + (c * cellWidth),
                origin.y - (r * cellHeight),
                0f
            );

            GameObject tile = Instantiate(wallPrefab, worldPos, 
                                          Quaternion.identity, transform);
            // Asignacion de Sprite segun tipo
            int spriteIndex = tileType - 1;
            if (spriteIndex >= 0 && spriteIndex < wallSprites.Length)
                tile.GetComponent<SpriteRenderer>().sprite = wallSprites[spriteIndex];
        }
    }
}
\end{lstlisting}
\end{frame}

% --- Mapeo de Sprites ---
\begin{frame}{Asignaci\'on Din\'amica de Sprites por Tipo de Bloque}
  Al instanciar el \texttt{wallPrefab}, se utiliza el valor num\'erico \texttt{tileType} le\'ido del mapa para configurar su representaci\'on visual:

  \vspace{0.3cm}
  \begin{table}
    \centering
    \begin{tabular}{|c|l|l|l|}
      \hline
      \textbf{Valor} & \textbf{Tipo de Bloque} & \textbf{\'Indice Sprite} & \textbf{Comportamiento} \\ \hline
      \texttt{0} & Fondo / Suelo & Ninguno & Espacio libre transitable \\ \hline
      \texttt{1} & Ladrillo & \texttt{wallSprites[0]} & Destructible por disparos \\ \hline
      \texttt{2} & Concreto / Acero & \texttt{wallSprites[1]} & Indestructible com\'un \\ \hline
      \texttt{3} & Pasto / Arbusto & \texttt{wallSprites[2]} & Oculta a los tanques \\ \hline
      \texttt{4} & Agua & \texttt{wallSprites[3]} & Bloquea tanques, no balas \\ \hline
    \end{tabular}
  \end{table}

  \vspace{0.3cm}
  \textbf{F\'ormula del \'indice:} $\text{spriteIndex} = \text{tileType} - 1$, validando siempre $0 \le \text{spriteIndex} < \text{wallSprites.Length}$.
\end{frame}

% --- Tabla Comparativa Global ---
\begin{frame}{Comparativa de Estructuras para Mapas en Unity}
  \begin{table}
    \centering
    \small
    \begin{tabular}{|l|l|l|}
      \hline
      \textbf{Estructura} & \textbf{Acceso / Recorrido} & \textbf{Caso de Uso en Unity} \\ \hline
      \texttt{T[]} (1D) & \texttt{arr[i]} & Lista de prefabs, sprites, armas \\ \hline
      \texttt{T[,]} (2D Rect.) & \texttt{arr[r, c]} & Tableros r\'igidos (Ajedrez, Damas) \\ \hline
      \texttt{T[][]} (Jagged) & \texttt{arr[r][c]} & Mapas din\'amicos le\'idos de \texttt{.txt} \\ \hline
      \texttt{TextAsset} & \texttt{.text} & Niveles integrados en el juego \\ \hline
      \texttt{System.IO.File} & \texttt{File.ReadAllLines} & Niveles externos / Mods de usuarios \\ \hline
    \end{tabular}
  \end{table}

  \vspace{0.3cm}
  \textbf{Recomendaci\'on:} Para niveles empaquetados en Unity, la combinaci\'on de \textbf{\texttt{TextAsset}} + \textbf{\texttt{int[][]}} ofrece el mejor equilibrio entre rendimiento, portabilidad y sencillez de c\'odigo.
\end{frame}

% --- Conclusiones ---
\begin{frame}{Conclusiones y Buenas Pr\'acticas}
  \begin{itemize}
    \item \textbf{Robustez con \texttt{int.TryParse}:} Evita que el juego se cierre si el archivo de texto contiene espacios adicionales o caracteres no num\'ericos.
    \item \textbf{Uso de \texttt{TextAsset}:} Asegura que los mapas funcionen en cualquier plataforma destino sin problemas de rutas relativas.
    \item \textbf{Cuidado con las coordenadas espaciales:} En una matriz la fila avanza hacia abajo ($+r$), lo que en el espacio mundial de Unity corresponde a restar altura ($-Y$).
    \item \textbf{Gesti\'on de memoria en escena:} Antes de cargar un nuevo nivel, destruir (\texttt{Destroy}) los bloques generados en la iteraci\'on anterior.
    \item \textbf{Flexibilidad de delimitadores:} Soportar comas, espacios y tabuladores hace que los archivos de texto sean f\'aciles de editar a mano por los dise\~nadores.
  \end{itemize}
\end{frame}

\end{document}
